% ==STYLE-BEGIN==
\documentclass[UTF8,11pt]{ctexart}
\usepackage{amsmath,amssymb,amsthm}
\usepackage[a4paper,margin=2.4cm,top=2.3cm,bottom=2.3cm]{geometry}
\usepackage{fancyhdr}
\usepackage{enumitem}
\usepackage{hyperref}
\hypersetup{hidelinks}

% ---- 版式 ----
\linespread{1.15}
\setlength{\parskip}{0.35em}
\setlength{\parindent}{1.6em}
\setlist{nosep,leftmargin=2em}

% ---- 页眉 ----
\pagestyle{fancy}
\fancyhf{}
\fancyhead[L]{\small\coursename}
\fancyhead[R]{\small\thepage}
\renewcommand{\headrulewidth}{0.4pt}

% ---- 定理环境：中文用正体，标题加粗，按小节共用一条流水号 ----
\newtheoremstyle{cnthm}{0.7em}{0.7em}{\normalfont}{}{\bfseries}{}{0.5em}{}
\theoremstyle{cnthm}
\newtheorem{thm}{定理}[section]
\newtheorem{lem}[thm]{引理}
\newtheorem{prop}[thm]{命题}
\newtheorem{cor}[thm]{推论}
\newtheorem{defn}[thm]{定义}
\newtheorem{rem}[thm]{注}
\newtheorem{exa}[thm]{例}
\newtheorem{conj}[thm]{猜想}
\renewcommand{\proofname}{\bfseries 证明}
\providecommand{\coursename}{}
% ==STYLE-END==
\renewcommand{\coursename}{解析几何}
\title{空间向量}
\date{}
\begin{document}
\maketitle

\section{\ 向量}

\begin{itemize}
\item \textbf{向量}\quad 有起点、有终点的有向线段，
一般用 $\vec a$ 或 $\overrightarrow{OP}$ 表示。向量两要素：长度，方向。
\item \textbf{零向量}\quad 长度为零的向量（起点与终点相同）。
\item \textbf{反向量}\quad 与给定向量长度相同、方向相反的向量。
\item 一个向量平移之后还是原来的向量。
\item \textbf{向量的加法}\quad 定义：平行四边形法则／三角形法则。
\end{itemize}

\subsubsection*{加法的性质}

\begin{enumerate}
\item 交换律\quad $\vec a+\vec b=\vec b+\vec a$；
\item 结合律\quad
$(\vec a+\vec b)+\vec c=\vec a+(\vec b+\vec c)$；

函数的复合是有结合律的！$f(gh)=(fg)h$。

\item 有零元\quad $\forall\vec a,\
\vec 0+\vec a=\vec a+\vec 0$；
\item 有负元\quad $\forall\vec a,\ \exists\vec b=-\vec a$，
s.t. $\vec a+\vec b=0$。
\end{enumerate}

代数系统 $=$ 非空集合 $+$ 运算 $+$ 法则。

\subsubsection*{向量的数乘}

$\lambda\in\mathbb R$，$\vec a$ 是空间向量，
$\lambda\cdot\vec a$ 是长度为 $|\lambda||\vec a|$、方向与 $\vec a$
相同（$\lambda>0$）或相反（$\lambda<0$）的向量。

显然 $0\cdot\vec 0=\vec 0$，$1\cdot\vec a=\vec a$，
$\lambda\cdot\vec 0=\vec 0$。

\subsubsection*{数乘的性质}

\begin{enumerate}
\item[5.] 有单位元\quad $\forall 1\cdot\vec a=\vec a$；
\item[6.] 分配律 1\quad
$\lambda(\vec a+\vec b)=\lambda\vec a+\lambda\vec b$；
\item[7.] 分配律 2\quad
$(\lambda+\mu)\vec a=\lambda\vec a+\mu\vec a$；
\item[8.] 结合律\quad
$(\lambda\mu)\vec a=\lambda(\mu\vec a)$；
\item[9.] 三角不等式\quad
$\bigl||\vec a|-|\vec b|\bigr|\leq|\vec a+\vec b|
\leq|\vec a|+|\vec b|$。
\end{enumerate}

线性空间的定义会用到这 8 条性质。

\subsubsection*{共线}

向量 $\vec a$ 与 $\vec b$ 共线 $\Longleftrightarrow$
$\vec a$ 与 $\vec b$ 方向相同或相反。

\textbf{定理一}\quad 若 $\vec a\neq\vec 0$，则 $\vec b$ 与
$\vec a$ 共线 $\Longleftrightarrow$ $\exists$ 唯一的
$t\in\mathbb R$，s.t. $\vec b=t\cdot\vec a$。

\end{document}

\end{document}
